\subsection{Topics 2 \& 3: CPU, Memory, Program Execution \& MIPS Assembly – Exercise Sheet [Solutions]}

\subsection*{Q\#1:}

Considering the simpler view of instructions, write the machine instructions representing each of the codes below:

\textbf{Memory locations/addresses:}
\begin{itemize}
  \item \texttt{num1 => 4040}
  \item \texttt{num2 => 5050}
  \item \texttt{Total => 6060}
\end{itemize}

\textbf{a.}
\begin{lstlisting}
int num1 = 5;
int num2 = 14;
int total = num1 + num2;
\end{lstlisting}

\textbf{Instruction:} \texttt{ADD 6060, 4040, 5050}

\textbf{b.}
\begin{lstlisting}
int num1 = 50;
int num2 = 14;
int total = num2 - num1;
\end{lstlisting}

\textbf{Instruction:} \texttt{SUB 6060, 5050, 4040}

\subsection*{Q\#2:}

Consider the below state of memory, where some values are already stored in different locations:

\begin{center}
\begin{tabular}{|c|c|c|c|c|c|c|c|c|}
\hline
\textbf{Value} & 50 & 70 & 14 & 9 & 7 & 44 & 20 \\
\hline
\textbf{Address} & 2024 & 2028 & 2032 & 2036 & 2040 & 2044 & 2048 & 2052 \\
\hline
\end{tabular}
\end{center}

What would be the result of executing the following instructions?

\begin{itemize}
  \item (a) \texttt{ADD 2028, 2024, 2048} \hfill [\textbf{94, Stored in 2028}]
  \item (b) \texttt{ADD 2036, 2024, 2036} \hfill [\textbf{64, Stored in 2036 and overwrites the 14}]
  \item (c) \texttt{MUL 2028, 2044, 2052} \hfill [\textbf{140, Stored in 2028 and overwrites the 94}]
  \item (d) \texttt{SUB 2040, 2048, 2044} \hfill [\textbf{37, Stored in 2040 and overwrites the 9}]
  \item (e) \texttt{XOR 2052, 2040, 2044} \hfill [\textbf{34, Stored in 2052}]
\end{itemize}

\textbf{The final memory state will look like:}

\begin{center}
\begin{tabular}{|c|c|c|c|c|c|c|c|c|}
\hline
\textbf{Value} & 50 & 140 & 70 & 64 & 37 & 7 & 44 & 34 \\
\hline
\textbf{Address} & 2024 & 2028 & 2032 & 2036 & 2040 & 2044 & 2048 & 2052 \\
\hline
\end{tabular}
\end{center}

\textbf{What would change if these instructions were to be executed in the following sequence:}

\begin{itemize}
  \item (a) \texttt{XOR 2052, 2040, 2044} \hfill [\textbf{14, Stored in 2052}]
  \item (b) \texttt{ADD 2036, 2024, 2036} \hfill [\textbf{64, Stored in 2036 and overwrites the 14}]
  \item (c) \texttt{ADD 2028, 2024, 2048} \hfill [\textbf{94, Stored in 2028}]
  \item (d) \texttt{SUB 2040, 2048, 2044} \hfill [\textbf{37, Stored in 2040 and overwrites the 9}]
  \item (e) \texttt{MUL 2028, 2044, 2052} \hfill [\textbf{98, Stored in 2028 and overwrites the 94}]
\end{itemize}

\textbf{When the instructions are executed in a different order, the resultant memory state will be different:}

\begin{center}
\begin{tabular}{|c|c|c|c|c|c|c|c|c|}
\hline
\textbf{Value} & 50 & 98 & 70 & 64 & 37 & 7 & 44 & 14 \\
\hline
\textbf{Address} & 2024 & 2028 & 2032 & 2036 & 2040 & 2044 & 2048 & 2052 \\
\hline
\end{tabular}
\end{center}

\subsection*{Q\#3: Instruction Cycle Description}

Using the figure below, describe the entire fetch / execute cycle detailing the result of each step in the cycle for the following machine instruction: \texttt{ADD 343, 546, 772}

\textbf{Instruction Fetch (IF):}
\begin{itemize}
  \item Execution begins by moving the instruction at the address given by the PC (600) from memory to the control unit.
  \item Bits of instruction are placed into the decoder circuit of the Control Unit.
  \item Once instruction is fetched, the PC can be readied for fetching the next instruction.
\end{itemize}

\textbf{Instruction Decode (ID):}
\begin{itemize}
  \item Decoder determines what operation the ALU will perform (ADD), and sets up the ALU.
  \item Decoder finds the memory addresses of the instruction's data (546, 772).
  \item Decoder finds the destination address for the Result Return step and places the address in the RR circuit (343).
\end{itemize}

\textbf{Data / Operand Fetch (DF/OF):}
\begin{itemize}
  \item The data values to be operated on are retrieved from memory (40, 23).
  \item Bits at specified memory locations are copied into locations in the ALU circuitry.
  \item Data values remain unchanged in the memory.
\end{itemize}

\textbf{Instruction Execution (EX):}
\begin{itemize}
  \item The addition circuit adds the two source operands (40, 23) together to produce their sum (63).
  \item Sum (63) is held in the ALU circuitry (accumulator).
\end{itemize}

\textbf{Result Return / Store (RR/ST):}
\begin{itemize}
  \item RR returns the result of EX (63) to the memory location specified by the destination address (343).
\end{itemize}

\subsection*{Q\#4: Pipeline Instruction Throughput}

A processor is operating at 32MHz. Each instruction takes a minimum of 6 cycles to execute. The processor has a six-stage pipeline. If a program starts execution at time 0, what is the theoretical maximum number of instructions that will have completed their execution at the end of 1 millisecond?

\textbf{Calculation:}
\begin{itemize}
  \item Speed = 32MHz = 32,000,000 cycles/second = 32,000 cycles/millisecond
  \item With pipelining:
    \begin{itemize}
      \item Number of clock cycles taken by the first instruction = 6 cycles
      \item After the first instruction has completely executed, one instruction comes out per cycle
      \item So, number of cycles taken by each remaining instruction = 1 cycle
    \end{itemize}
  \item Number of executed instructions = First Instruction (1) + remaining instructions (32,000 cycles/ms – 6 cycles) = $1 + 31,994 = \boxed{31,995\ \text{instructions/ms}}$
\end{itemize}

\textit{In practice, there are factors that mean this theoretical maximum is never reached.}

\subsection*{Q\#5: Loop Result in MIPS Assembly}

\textbf{What is the value stored in \texttt{z} after the execution of the following MIPS assembly code, if \texttt{x=5}?}

\begin{lstlisting}
li $t0, 1
li $t1, 2
lw $t2, &x
L1: bgt $t1, $t2, L2
    mul $t0, $t0, $t1
    addi $t1, $t1, 1
    j L1
L2: sw $t0, &z
\end{lstlisting}

\textbf{Explanation:}  
The \texttt{bgt} instruction compares two registers and jumps to \texttt{L2} if the content of \texttt{\$t1} is greater than the content of \texttt{\$t2}. Equivalent Java code is:

\begin{lstlisting}[language=Java]
z = 1;
for (i = 2; i <= x; i++) {
    z = z * i;
}
\end{lstlisting}

Knowing that \texttt{x = 5} and the index \texttt{i} starts from 2, the result will be:  
\texttt{z = 1 * 2 * 3 * 4 * 5 = 120}

\subsection*{Q\#6: Expression to MIPS Translation}

\textbf{Expression:}  
\[
C = 12 \times A + 4 \times (B - 5)
\]

\textbf{Breakdown:}
\begin{align*}
t0 &= 12 \times A \\
t1 &= B - 5 \\
t2 &= 4 \times t1 \\
C &= t0 + t2
\end{align*}

\textbf{MIPS Assembly Code and Execution Results:}

\begin{lstlisting}
lw $s0, &A        # $s0 = 10
lw $s1, &B        # $s1 = 9
mul $t0, $s0, 12  # $t0 = 120
sub $t1, $s1, 5   # $t1 = 4
mul $t2, $t1, 4   # $t2 = 16
add $s2, $t0, $t2 # $s2 = 136
sw $s2, &C        # store result
\end{lstlisting}

\textbf{Result:} \texttt{C = 136}

